% ============================================================================
% ARBEITSBLATT - Physik Klasse 6: Weg-Zeit-Diagramme
% ============================================================================
% Begleitend zu LP-003, mit Merkkasten + Übungsaufgaben
% ============================================================================

\documentclass[11pt,a4paper]{article}

\usepackage[utf8]{inputenc}
\usepackage[T1]{fontenc}
\usepackage[ngerman]{babel}
\usepackage{geometry}
\usepackage{helvet}
\usepackage{tikz}
\usepackage{enumitem}
\usepackage{fancyhdr}
\usepackage{xcolor}
\usepackage{tcolorbox}
\usepackage{amssymb}

\renewcommand{\familydefault}{\sfdefault}

\geometry{left=1.8cm, right=1.8cm, top=1.8cm, bottom=1.5cm}

% Farben
\definecolor{physik}{HTML}{2c5aa0}
\definecolor{lightblue}{HTML}{e8f0fa}
\definecolor{gridcolor}{HTML}{c0d0e8}
\definecolor{gridlight}{HTML}{e8eef5}

% Kopfzeile (schwarz auf weiß)
\pagestyle{fancy}
\fancyhf{}
\fancyhead[L]{\small\textbf{Physik Klasse 6} | Weg-Zeit-Diagramme}
\fancyhead[R]{\small Arbeitsblatt}
\renewcommand{\headrulewidth}{0.4pt}

% Lücke
\newcommand{\luecke}[1]{\underline{\hspace{#1}}}

\tcbuselibrary{skins}

% Merkkasten (weißer Hintergrund, farbiger Rahmen, schwarze Überschrift)
\newtcolorbox{merkkasten}[1][]{
    colback=white,
    colframe=physik,
    coltitle=black,
    colbacktitle=white,
    fonttitle=\bfseries\small,
    boxrule=1pt,
    arc=2pt,
    left=8pt, right=8pt, top=6pt, bottom=6pt,
    #1
}

\begin{document}

% ============================================================================
% TITEL
% ============================================================================
\begin{center}
    {\Large\bfseries Weg-Zeit-Diagramme verstehen}\\[0.2em]
    {\small Name: \luecke{5cm} \hfill Datum: \luecke{3cm}}
\end{center}
\vspace{0.2em}
\hrule
\vspace{0.8em}

% ============================================================================
% KASTEN 1: Achsenbeschriftung
% ============================================================================
\begin{merkkasten}[title=Kasten 1: Aufbau eines Weg-Zeit-Diagramms]
\begin{minipage}[t]{0.55\textwidth}
\vspace{0pt}
\begin{itemize}[leftmargin=*, itemsep=3pt, topsep=0pt]
    \item Die \textbf{waagerechte Achse} (x-Achse) zeigt die \luecke{1.5cm}~$t$ in \luecke{1cm}.
    \item Die \textbf{senkrechte Achse} (y-Achse) zeigt den \luecke{1.5cm}~$s$ in \luecke{1cm}.
\end{itemize}
\end{minipage}
\hfill
\begin{minipage}[t]{0.40\textwidth}
\vspace{0pt}
\begin{tikzpicture}[scale=0.4]
    \draw[gridcolor, very thin] (0,0) grid (6,4);
    \draw[->, thick] (0,0) -- (6.5,0) node[right] {\small $t$ in s};
    \draw[->, thick] (0,0) -- (0,4.5) node[above] {\small $s$ in m};
    \node[below] at (3,-0.5) {\tiny waagerecht};
    \node[left, rotate=90] at (-0.8,2) {\tiny senkrecht};
\end{tikzpicture}
\end{minipage}
\end{merkkasten}

\vspace{0.6em}

% ============================================================================
% KASTEN 2: Interpretation schnell/langsam
% ============================================================================
\begin{merkkasten}[title=Kasten 2: Schnell oder langsam?]
\begin{minipage}[t]{0.55\textwidth}
\vspace{0pt}
Die \textbf{Steigung} der Geraden zeigt die Geschwindigkeit:
\begin{itemize}[leftmargin=*, itemsep=2pt, topsep=3pt]
    \item \textbf{Steile} Gerade = \luecke{2cm} Bewegung
    \item \textbf{Flache} Gerade = \luecke{2cm} Bewegung
\end{itemize}
\end{minipage}
\hfill
\begin{minipage}[t]{0.40\textwidth}
\vspace{0pt}
\begin{tikzpicture}[scale=0.38]
    \draw[gridcolor, very thin] (0,0) grid (6,4);
    \draw[->, thick] (0,0) -- (6.5,0) node[right] {\tiny $t$};
    \draw[->, thick] (0,0) -- (0,4.5) node[above] {\tiny $s$};
    \draw[red!70!black, very thick] (0,0) -- (3,3.5);
    \node[red!70!black, right] at (3.2,3.5) {\tiny schnell};
    \draw[physik, very thick] (0,0) -- (6,2);
    \node[physik, right] at (6.2,2) {\tiny langsam};
\end{tikzpicture}
\end{minipage}
\end{merkkasten}

\vspace{0.6em}

% ============================================================================
% KASTEN 3: Vier Bewegungsarten
% ============================================================================
\begin{merkkasten}[title=Kasten 3: Vier Bewegungsarten im Diagramm]
\begin{center}
\begin{tikzpicture}[scale=0.38]
    % Diagramm 1: Schnell
    \begin{scope}[shift={(0,0)}]
        \draw[gridcolor, very thin] (0,0) grid (4,3);
        \draw[->, thick] (0,0) -- (4.3,0);
        \draw[->, thick] (0,0) -- (0,3.3);
        \draw[physik, very thick] (0,0) -- (4,3);
        \node[below] at (2,-0.8) {\tiny \textbf{steil}};
        \node[below] at (2,-2.2) {\tiny = \luecke{1.4cm}};
    \end{scope}
    % Diagramm 2: Langsam
    \begin{scope}[shift={(7.5,0)}]
        \draw[gridcolor, very thin] (0,0) grid (4,3);
        \draw[->, thick] (0,0) -- (4.3,0);
        \draw[->, thick] (0,0) -- (0,3.3);
        \draw[physik, very thick] (0,0) -- (4,1.2);
        \node[below] at (2,-0.8) {\tiny \textbf{flach}};
        \node[below] at (2,-2.2) {\tiny = \luecke{1.4cm}};
    \end{scope}
    % Diagramm 3: Stillstand
    \begin{scope}[shift={(15,0)}]
        \draw[gridcolor, very thin] (0,0) grid (4,3);
        \draw[->, thick] (0,0) -- (4.3,0);
        \draw[->, thick] (0,0) -- (0,3.3);
        \draw[physik, very thick] (0,1.5) -- (4,1.5);
        \node[below] at (2,-0.8) {\tiny \textbf{waagerecht}};
        \node[below] at (2,-2.2) {\tiny = \luecke{1.4cm}};
    \end{scope}
    % Diagramm 4: Ungleichförmig
    \begin{scope}[shift={(22.5,0)}]
        \draw[gridcolor, very thin] (0,0) grid (4,3);
        \draw[->, thick] (0,0) -- (4.3,0);
        \draw[->, thick] (0,0) -- (0,3.3);
        \draw[physik, very thick] (0,0) .. controls (1.5,0.5) and (2.5,2) .. (4,3);
        \node[below] at (2,-0.8) {\tiny \textbf{Kurve}};
        \node[below] at (2,-2.2) {\tiny = \luecke{1.6cm}};
    \end{scope}
\end{tikzpicture}
\end{center}
\end{merkkasten}

\vspace{1em}

% ============================================================================
% AUFGABE 1: Ablesen üben
% ============================================================================
\noindent\textbf{\textcolor{physik}{Aufgabe 1: Werte ablesen}}

\vspace{0.3em}
\begin{minipage}[t]{0.58\textwidth}
\noindent Ein Radfahrer fährt gleichförmig.

\vspace{0.3em}
\begin{tikzpicture}[scale=0.7]
    \draw[gridlight, very thin] (0,0) grid[step=0.5] (8,5);
    \draw[gridcolor, thin] (0,0) grid (8,5);
    \draw[->, thick] (0,0) -- (8.5,0) node[right] {\small $t$ in s};
    \draw[->, thick] (0,0) -- (0,5.5) node[above] {\small $s$ in m};
    \foreach \x/\label in {0/0, 2/2, 4/4, 6/6, 8/8} {
        \draw[thick] (\x,-0.1) -- (\x,0.1);
        \node[below] at (\x,-0.2) {\small \label};
    }
    \foreach \y/\label in {0/0, 1/10, 2/20, 3/30, 4/40, 5/50} {
        \draw[thick] (-0.1,\y) -- (0.1,\y);
        \node[left] at (-0.2,\y) {\small \label};
    }
    \draw[physik, very thick] (0,0) -- (8,4);
    \foreach \x/\y in {0/0, 2/1, 4/2, 6/3, 8/4} {
        \filldraw[physik] (\x,\y) circle (2.5pt);
    }
\end{tikzpicture}
\end{minipage}
\hfill
\begin{minipage}[t]{0.38\textwidth}
\vspace{2em}
\noindent\textbf{Lies ab:}
\begin{enumerate}[label=\alph*), leftmargin=1.5em, itemsep=0.6em]
    \item Weg nach $t = 4$ s? \quad $s =$ \luecke{1.5cm}
    \item Zeit bei $s = 30$ m? \quad $t =$ \luecke{1.5cm}
    \item Weg nach $t = 8$ s? \quad $s =$ \luecke{1.5cm}
\end{enumerate}
\end{minipage}

\newpage

% ============================================================================
% AUFGABE 2: Diagramm zeichnen
% ============================================================================
\noindent\textbf{\textcolor{physik}{Aufgabe 2: Ein Diagramm zeichnen}}

\vspace{0.3em}
\noindent Ein Jogger läuft gleichförmig. Trage die Punkte ein und verbinde sie.

\vspace{0.5em}
\begin{center}
\begin{tabular}{|c|c|c|c|c|c|}
\hline
$t$ in s & 0 & 4 & 8 & 12 & 16 \\
\hline
$s$ in m & 0 & 20 & 40 & 60 & 80 \\
\hline
\end{tabular}
\end{center}

\vspace{0.5em}
\begin{center}
\begin{tikzpicture}[scale=0.68]
    \def\xmax{16}
    \def\ymax{8}
    \foreach \x in {0,0.5,...,\xmax} {
        \draw[gridlight, very thin] (\x,0) -- (\x,\ymax);
    }
    \foreach \y in {0,0.5,...,\ymax} {
        \draw[gridlight, very thin] (0,\y) -- (\xmax,\y);
    }
    \foreach \x in {0,1,...,\xmax} {
        \draw[gridcolor, thin] (\x,0) -- (\x,\ymax);
    }
    \foreach \y in {0,1,...,\ymax} {
        \draw[gridcolor, thin] (0,\y) -- (\xmax,\y);
    }
    \draw[->, thick] (0,0) -- (\xmax+0.6,0) node[right] {\small $t$ in s};
    \draw[->, thick] (0,0) -- (0,\ymax+0.6) node[above] {\small $s$ in m};
    \foreach \x/\label in {0/0, 4/4, 8/8, 12/12, 16/16} {
        \draw[thick] (\x,-0.1) -- (\x,0.1);
        \node[below] at (\x,-0.25) {\small \label};
    }
    \foreach \y/\label in {0/0, 2/20, 4/40, 6/60, 8/80} {
        \draw[thick] (-0.1,\y) -- (0.1,\y);
        \node[left] at (-0.25,\y) {\small \label};
    }
\end{tikzpicture}
\end{center}

\vspace{1.2em}

% ============================================================================
% AUFGABE 3 und 4 (vormals Erweiterungsaufgaben)
% ============================================================================

\begin{minipage}[t]{0.48\textwidth}
\noindent\textbf{\textcolor{physik}{Aufgabe 3: Ohne Zahlen}} -- Was passiert?

\vspace{0.2em}
\begin{tikzpicture}[scale=0.35]
    \draw[gridcolor, very thin] (0,0) grid (8,4);
    \draw[->, thick] (0,0) -- (8.5,0) node[right] {\tiny $t$};
    \draw[->, thick] (0,0) -- (0,4.5) node[above] {\tiny $s$};
    \draw[physik, very thick] (0,0) -- (2.4,2.4);
    \draw[physik, very thick] (2.4,2.4) -- (4.8,2.4);
    \draw[physik, very thick] (4.8,2.4) -- (8,3.6);
\end{tikzpicture}

\vspace{0.2em}
{\footnotesize
1: \luecke{2.5cm}\\[0.3em]
2: \luecke{2.5cm}\\[0.3em]
3: \luecke{2.5cm}
}
\end{minipage}
\hfill
\begin{minipage}[t]{0.48\textwidth}
\noindent\textbf{\textcolor{physik}{Aufgabe 4: Stillstand oder gleichförmig?}}

\vspace{0.2em}
\begin{tikzpicture}[scale=0.35]
    \draw[gridcolor, very thin] (0,0) grid (8,4);
    \draw[->, thick] (0,0) -- (8.5,0) node[right] {\tiny $t$};
    \draw[->, thick] (0,0) -- (0,4.5) node[above] {\tiny $s$};
    \draw[physik, very thick] (0,0) -- (8,0);
    \node[physik, right] at (8.2,0) {\tiny A};
    \draw[red!70!black, very thick] (0,0) -- (8,3);
    \node[red!70!black, right] at (8.2,3) {\tiny B};
\end{tikzpicture}

\vspace{0.2em}
{\footnotesize
A: $\square$ gleichförmig \quad $\square$ Stillstand\\[0.3em]
B: $\square$ gleichförmig \quad $\square$ Stillstand
}
\end{minipage}

\end{document}
